\begin{examplebox}
\textbf{\textcolor{CExample}{例 1（基础）}}

\textbf{\textcolor{CExample}{题目：}} 椭圆方程为 \(\frac{x^{2}}{16}+\frac{y^{2}}{9}=1\)，将其绕 \(x\) 轴旋转一周，求所得旋转体的体积。

\textbf{\textcolor{CExample}{解题步骤：}}
\begin{description}[style=nextline, leftmargin=2.8em, labelsep=0.5em, itemsep=0.45em]
    \item[\textbf{第一步（识别半轴）：}] 由标准方程直接读出 \(a^{2}=16\)，\(b^{2}=9\)，故 \(a=4\)，\(b=3\)。
    \par\textbf{理由：} 标准形式 \(\frac{x^{2}}{a^{2}}+\frac{y^{2}}{b^{2}}=1\) 中分母即为半轴平方。
    \item[\textbf{第二步（套用公式）：}] 旋转椭球体积公式为 \(V=\frac{4}{3}\pi a b^{2}\)，代入 \(a=4\)，\(b=3\)：
    \[
    V=\frac{4}{3}\pi \times 4 \times 3^{2}.
    \]
    \par\textbf{理由：} 绕 \(x\) 轴旋转时适用该公式。
    \item[\textbf{第三步（计算结果）：}] 计算数值：
    \[
    V=\frac{4}{3}\pi \times 4 \times 9 = \frac{4}{3}\pi \times 36 = 48\pi.
    \]
    \par\textbf{理由：} 先算指数，再约分相乘。
\end{description}

\textbf{\textcolor{CExample}{关键结论：}} \textbf{所求体积为 \(48\pi\)（体积单位）。}

\medskip
\textbf{\textcolor{CExample}{例 2（稍变形）}}

\textbf{\textcolor{CExample}{题目：}} 已知椭圆方程为 \(9x^{2}+25y^{2}=225\)，将该椭圆绕 \(x\) 轴旋转，求旋转椭球体积。

\textbf{\textcolor{CExample}{解题步骤：}}
\begin{description}[style=nextline, leftmargin=2.8em, labelsep=0.5em, itemsep=0.45em]
    \item[\textbf{第一步（化为标准式）：}] 方程两边同除以 \(225\)，得 \(\frac{x^{2}}{25}+\frac{y^{2}}{9}=1\)。
    \par\textbf{理由：} 必须化为 \(\frac{x^{2}}{a^{2}}+\frac{y^{2}}{b^{2}}=1\) 的标准形式才能准确识别 \(a,b\)。
    \item[\textbf{第二步（提取半轴）：}] 由此知 \(a^{2}=25\)，\(b^{2}=9\)，即 \(a=5\)，\(b=3\)。
    \par\textbf{理由：} \(a,b>0\)。
    \item[\textbf{第三步（代入体积公式）：}] 
    \[
    V=\frac{4}{3}\pi a b^{2} = \frac{4}{3}\pi \times 5 \times 3^{2} = \frac{4}{3}\pi \times 5 \times 9 = 60\pi.
    \]
    \par\textbf{理由：} 套用公式直接计算。
\end{description}

\textbf{\textcolor{CExample}{关键结论：}} \textbf{旋转椭球体积为 \(60\pi\)。}

\medskip
\textbf{\textcolor{CExample}{例 3（综合应用）}}

\textbf{\textcolor{CExample}{题目：}} 一椭圆绕 \(x\) 轴旋转一周所得椭球体积为 \(\frac{16\pi}{3}\)，已知椭圆在 \(x\) 轴上的半轴长为 \(2\)，求椭圆标准方程。

\textbf{\textcolor{CExample}{解题步骤：}}
\begin{description}[style=nextline, leftmargin=2.8em, labelsep=0.5em, itemsep=0.45em]
    \item[\textbf{第一步（设未知数列方程）：}] 设椭圆标准方程为 \(\frac{x^{2}}{a^{2}}+\frac{y^{2}}{b^{2}}=1\)，已知 \(a=2\)，体积 \(V=\frac{16\pi}{3}\)。由体积公式 
    \[
    V=\frac{4}{3}\pi a b^{2},
    \]
    代入得 \(\frac{4}{3}\pi \times 2 \times b^{2} = \frac{16\pi}{3}\)。
    \par\textbf{理由：} 逆用公式建立关于 \(b^{2}\) 的方程。
    \item[\textbf{第二步（化简求解）：}] 两边消去 \(\pi\) 并化简：
    \[
    \frac{8}{3} b^{2} = \frac{16}{3} \quad\Longrightarrow\quad b^{2}=2.
    \]
    \par\textbf{理由：} 等式两边同乘 \(\frac{3}{8}\)，或交叉相乘。
    \item[\textbf{第三步（写出方程）：}] 将 \(a^{2}=4\)，\(b^{2}=2\) 代回标准形式：
    \[
    \frac{x^{2}}{4}+\frac{y^{2}}{2}=1.
    \]
    \par\textbf{理由：} 半轴平方 \(a^{2}=2^{2}=4\)，\(b^{2}=2\)。
\end{description}

\textbf{\textcolor{CExample}{关键结论：}} \textbf{所求椭圆方程为 \(\frac{x^{2}}{4}+\frac{y^{2}}{2}=1\)。}

\medskip
\textbf{\textcolor{CExample}{例 4（同一题的两种新增验证法）}}

\textbf{\textcolor{CExample}{题目：}} 仍取例 1 的椭圆 \(\frac{x^{2}}{16}+\frac{y^{2}}{9}=1\)，将其绕 \(x\) 轴旋转。除原有圆盘法外，分别用沿 \(y\) 方向的椭圆截面法、球体伸缩法验证体积。

\textbf{\textcolor{CExample}{解法一（沿 y 方向的椭圆截面积分）：}}
\begin{description}[style=nextline, leftmargin=2.8em, labelsep=0.5em, itemsep=0.45em]
    \item[\textbf{第一步（写出截面）：}] 三维方程是 \(\frac{x^{2}}{16}+\frac{y^{2}+z^{2}}{9}=1\)。在 \(y=c\) 处截面边界为 \(\frac{x^{2}}{16}+\frac{z^{2}}{9}=1-\frac{c^{2}}{9}\)，两条半轴为 \(4\sqrt{1-c^{2}/9}\) 与 \(3\sqrt{1-c^{2}/9}\)。
    \item[\textbf{第二步（椭圆面积积分）：}]
    \[
    \begin{aligned}
    S(y)&=12\pi\left(1-\frac{y^{2}}{9}\right),\\
    V&=12\pi\int_{-3}^{3}\left(1-\frac{y^{2}}{9}\right)dy
      =12\pi\left[y-\frac{y^{3}}{27}\right]_{-3}^{3}
      =\boxed{48\pi}.
    \end{aligned}
    \]
\end{description}

\textbf{\textcolor{CExample}{解法二（半径为 3 的球沿 x 方向伸缩）：}}
半径为 \(3\) 的球体积为 \(\frac43\pi\cdot3^{3}=36\pi\)。把球的 \(x\) 方向半轴从 \(3\) 伸缩到 \(4\)，倍数为 \(4/3\)。变换 \(x=\frac43u\) 把球面方程 \(\frac{u^{2}+y^{2}+z^{2}}{9}=1\) 变成 \(\frac{x^{2}}{16}+\frac{y^{2}+z^{2}}{9}=1\)。故
\[
V=\frac43\cdot36\pi=\boxed{48\pi}.
\]
\textbf{\textcolor{CExample}{关键结论：}} \textbf{同一椭球换切片方向，或者由球伸缩而来，体积都为 \(48\pi\)。}

\end{examplebox}
